\documentclass{article} % For LaTeX2e
\usepackage{preprint}  % options: [anonymous], [monochrome], [wide]
\usepackage{microtype}

% Optional math commands from https://github.com/goodfeli/dlbook_notation.
%%%%% NEW MATH DEFINITIONS %%%%%

\usepackage{amsmath,amsfonts,bm}

% Mark sections of captions for referring to divisions of figures
\newcommand{\figleft}{{\em (Left)}}
\newcommand{\figcenter}{{\em (Center)}}
\newcommand{\figright}{{\em (Right)}}
\newcommand{\figtop}{{\em (Top)}}
\newcommand{\figbottom}{{\em (Bottom)}}
\newcommand{\captiona}{{\em (a)}}
\newcommand{\captionb}{{\em (b)}}
\newcommand{\captionc}{{\em (c)}}
\newcommand{\captiond}{{\em (d)}}

% Highlight a newly defined term
\newcommand{\newterm}[1]{{\bf #1}}


% Figure reference, lower-case.
\def\figref#1{figure~\ref{#1}}
% Figure reference, capital. For start of sentence
\def\Figref#1{Figure~\ref{#1}}
\def\twofigref#1#2{figures \ref{#1} and \ref{#2}}
\def\quadfigref#1#2#3#4{figures \ref{#1}, \ref{#2}, \ref{#3} and \ref{#4}}
% Section reference, lower-case.
\def\secref#1{section~\ref{#1}}
% Section reference, capital.
\def\Secref#1{Section~\ref{#1}}
% Reference to two sections.
\def\twosecrefs#1#2{sections \ref{#1} and \ref{#2}}
% Reference to three sections.
\def\secrefs#1#2#3{sections \ref{#1}, \ref{#2} and \ref{#3}}
% Reference to an equation, lower-case.
\def\eqref#1{equation~\ref{#1}}
% Reference to an equation, upper case
\def\Eqref#1{Equation~\ref{#1}}
% A raw reference to an equation---avoid using if possible
\def\plaineqref#1{\ref{#1}}
% Reference to a chapter, lower-case.
\def\chapref#1{chapter~\ref{#1}}
% Reference to an equation, upper case.
\def\Chapref#1{Chapter~\ref{#1}}
% Reference to a range of chapters
\def\rangechapref#1#2{chapters\ref{#1}--\ref{#2}}
% Reference to an algorithm, lower-case.
\def\algref#1{algorithm~\ref{#1}}
% Reference to an algorithm, upper case.
\def\Algref#1{Algorithm~\ref{#1}}
\def\twoalgref#1#2{algorithms \ref{#1} and \ref{#2}}
\def\Twoalgref#1#2{Algorithms \ref{#1} and \ref{#2}}
% Reference to a part, lower case
\def\partref#1{part~\ref{#1}}
% Reference to a part, upper case
\def\Partref#1{Part~\ref{#1}}
\def\twopartref#1#2{parts \ref{#1} and \ref{#2}}

\def\ceil#1{\lceil #1 \rceil}
\def\floor#1{\lfloor #1 \rfloor}
\def\1{\bm{1}}
\newcommand{\train}{\mathcal{D}}
\newcommand{\valid}{\mathcal{D_{\mathrm{valid}}}}
\newcommand{\test}{\mathcal{D_{\mathrm{test}}}}

\def\eps{{\epsilon}}


% Random variables
\def\reta{{\textnormal{$\eta$}}}
\def\ra{{\textnormal{a}}}
\def\rb{{\textnormal{b}}}
\def\rc{{\textnormal{c}}}
\def\rd{{\textnormal{d}}}
\def\re{{\textnormal{e}}}
\def\rf{{\textnormal{f}}}
\def\rg{{\textnormal{g}}}
\def\rh{{\textnormal{h}}}
\def\ri{{\textnormal{i}}}
\def\rj{{\textnormal{j}}}
\def\rk{{\textnormal{k}}}
\def\rl{{\textnormal{l}}}
% rm is already a command, just don't name any random variables m
\def\rn{{\textnormal{n}}}
\def\ro{{\textnormal{o}}}
\def\rp{{\textnormal{p}}}
\def\rq{{\textnormal{q}}}
\def\rr{{\textnormal{r}}}
\def\rs{{\textnormal{s}}}
\def\rt{{\textnormal{t}}}
\def\ru{{\textnormal{u}}}
\def\rv{{\textnormal{v}}}
\def\rw{{\textnormal{w}}}
\def\rx{{\textnormal{x}}}
\def\ry{{\textnormal{y}}}
\def\rz{{\textnormal{z}}}

% Random vectors
\def\rvepsilon{{\mathbf{\epsilon}}}
\def\rvtheta{{\mathbf{\theta}}}
\def\rva{{\mathbf{a}}}
\def\rvb{{\mathbf{b}}}
\def\rvc{{\mathbf{c}}}
\def\rvd{{\mathbf{d}}}
\def\rve{{\mathbf{e}}}
\def\rvf{{\mathbf{f}}}
\def\rvg{{\mathbf{g}}}
\def\rvh{{\mathbf{h}}}
\def\rvu{{\mathbf{i}}}
\def\rvj{{\mathbf{j}}}
\def\rvk{{\mathbf{k}}}
\def\rvl{{\mathbf{l}}}
\def\rvm{{\mathbf{m}}}
\def\rvn{{\mathbf{n}}}
\def\rvo{{\mathbf{o}}}
\def\rvp{{\mathbf{p}}}
\def\rvq{{\mathbf{q}}}
\def\rvr{{\mathbf{r}}}
\def\rvs{{\mathbf{s}}}
\def\rvt{{\mathbf{t}}}
\def\rvu{{\mathbf{u}}}
\def\rvv{{\mathbf{v}}}
\def\rvw{{\mathbf{w}}}
\def\rvx{{\mathbf{x}}}
\def\rvy{{\mathbf{y}}}
\def\rvz{{\mathbf{z}}}

% Elements of random vectors
\def\erva{{\textnormal{a}}}
\def\ervb{{\textnormal{b}}}
\def\ervc{{\textnormal{c}}}
\def\ervd{{\textnormal{d}}}
\def\erve{{\textnormal{e}}}
\def\ervf{{\textnormal{f}}}
\def\ervg{{\textnormal{g}}}
\def\ervh{{\textnormal{h}}}
\def\ervi{{\textnormal{i}}}
\def\ervj{{\textnormal{j}}}
\def\ervk{{\textnormal{k}}}
\def\ervl{{\textnormal{l}}}
\def\ervm{{\textnormal{m}}}
\def\ervn{{\textnormal{n}}}
\def\ervo{{\textnormal{o}}}
\def\ervp{{\textnormal{p}}}
\def\ervq{{\textnormal{q}}}
\def\ervr{{\textnormal{r}}}
\def\ervs{{\textnormal{s}}}
\def\ervt{{\textnormal{t}}}
\def\ervu{{\textnormal{u}}}
\def\ervv{{\textnormal{v}}}
\def\ervw{{\textnormal{w}}}
\def\ervx{{\textnormal{x}}}
\def\ervy{{\textnormal{y}}}
\def\ervz{{\textnormal{z}}}

% Random matrices
\def\rmA{{\mathbf{A}}}
\def\rmB{{\mathbf{B}}}
\def\rmC{{\mathbf{C}}}
\def\rmD{{\mathbf{D}}}
\def\rmE{{\mathbf{E}}}
\def\rmF{{\mathbf{F}}}
\def\rmG{{\mathbf{G}}}
\def\rmH{{\mathbf{H}}}
\def\rmI{{\mathbf{I}}}
\def\rmJ{{\mathbf{J}}}
\def\rmK{{\mathbf{K}}}
\def\rmL{{\mathbf{L}}}
\def\rmM{{\mathbf{M}}}
\def\rmN{{\mathbf{N}}}
\def\rmO{{\mathbf{O}}}
\def\rmP{{\mathbf{P}}}
\def\rmQ{{\mathbf{Q}}}
\def\rmR{{\mathbf{R}}}
\def\rmS{{\mathbf{S}}}
\def\rmT{{\mathbf{T}}}
\def\rmU{{\mathbf{U}}}
\def\rmV{{\mathbf{V}}}
\def\rmW{{\mathbf{W}}}
\def\rmX{{\mathbf{X}}}
\def\rmY{{\mathbf{Y}}}
\def\rmZ{{\mathbf{Z}}}

% Elements of random matrices
\def\ermA{{\textnormal{A}}}
\def\ermB{{\textnormal{B}}}
\def\ermC{{\textnormal{C}}}
\def\ermD{{\textnormal{D}}}
\def\ermE{{\textnormal{E}}}
\def\ermF{{\textnormal{F}}}
\def\ermG{{\textnormal{G}}}
\def\ermH{{\textnormal{H}}}
\def\ermI{{\textnormal{I}}}
\def\ermJ{{\textnormal{J}}}
\def\ermK{{\textnormal{K}}}
\def\ermL{{\textnormal{L}}}
\def\ermM{{\textnormal{M}}}
\def\ermN{{\textnormal{N}}}
\def\ermO{{\textnormal{O}}}
\def\ermP{{\textnormal{P}}}
\def\ermQ{{\textnormal{Q}}}
\def\ermR{{\textnormal{R}}}
\def\ermS{{\textnormal{S}}}
\def\ermT{{\textnormal{T}}}
\def\ermU{{\textnormal{U}}}
\def\ermV{{\textnormal{V}}}
\def\ermW{{\textnormal{W}}}
\def\ermX{{\textnormal{X}}}
\def\ermY{{\textnormal{Y}}}
\def\ermZ{{\textnormal{Z}}}

% Vectors
\def\vzero{{\bm{0}}}
\def\vone{{\bm{1}}}
\def\vmu{{\bm{\mu}}}
\def\vtheta{{\bm{\theta}}}
\def\va{{\bm{a}}}
\def\vb{{\bm{b}}}
\def\vc{{\bm{c}}}
\def\vd{{\bm{d}}}
\def\ve{{\bm{e}}}
\def\vf{{\bm{f}}}
\def\vg{{\bm{g}}}
\def\vh{{\bm{h}}}
\def\vi{{\bm{i}}}
\def\vj{{\bm{j}}}
\def\vk{{\bm{k}}}
\def\vl{{\bm{l}}}
\def\vm{{\bm{m}}}
\def\vn{{\bm{n}}}
\def\vo{{\bm{o}}}
\def\vp{{\bm{p}}}
\def\vq{{\bm{q}}}
\def\vr{{\bm{r}}}
\def\vs{{\bm{s}}}
\def\vt{{\bm{t}}}
\def\vu{{\bm{u}}}
\def\vv{{\bm{v}}}
\def\vw{{\bm{w}}}
\def\vx{{\bm{x}}}
\def\vy{{\bm{y}}}
\def\vz{{\bm{z}}}

% Elements of vectors
\def\evalpha{{\alpha}}
\def\evbeta{{\beta}}
\def\evepsilon{{\epsilon}}
\def\evlambda{{\lambda}}
\def\evomega{{\omega}}
\def\evmu{{\mu}}
\def\evpsi{{\psi}}
\def\evsigma{{\sigma}}
\def\evtheta{{\theta}}
\def\eva{{a}}
\def\evb{{b}}
\def\evc{{c}}
\def\evd{{d}}
\def\eve{{e}}
\def\evf{{f}}
\def\evg{{g}}
\def\evh{{h}}
\def\evi{{i}}
\def\evj{{j}}
\def\evk{{k}}
\def\evl{{l}}
\def\evm{{m}}
\def\evn{{n}}
\def\evo{{o}}
\def\evp{{p}}
\def\evq{{q}}
\def\evr{{r}}
\def\evs{{s}}
\def\evt{{t}}
\def\evu{{u}}
\def\evv{{v}}
\def\evw{{w}}
\def\evx{{x}}
\def\evy{{y}}
\def\evz{{z}}

% Matrix
\def\mA{{\bm{A}}}
\def\mB{{\bm{B}}}
\def\mC{{\bm{C}}}
\def\mD{{\bm{D}}}
\def\mE{{\bm{E}}}
\def\mF{{\bm{F}}}
\def\mG{{\bm{G}}}
\def\mH{{\bm{H}}}
\def\mI{{\bm{I}}}
\def\mJ{{\bm{J}}}
\def\mK{{\bm{K}}}
\def\mL{{\bm{L}}}
\def\mM{{\bm{M}}}
\def\mN{{\bm{N}}}
\def\mO{{\bm{O}}}
\def\mP{{\bm{P}}}
\def\mQ{{\bm{Q}}}
\def\mR{{\bm{R}}}
\def\mS{{\bm{S}}}
\def\mT{{\bm{T}}}
\def\mU{{\bm{U}}}
\def\mV{{\bm{V}}}
\def\mW{{\bm{W}}}
\def\mX{{\bm{X}}}
\def\mY{{\bm{Y}}}
\def\mZ{{\bm{Z}}}
\def\mBeta{{\bm{\beta}}}
\def\mPhi{{\bm{\Phi}}}
\def\mLambda{{\bm{\Lambda}}}
\def\mSigma{{\bm{\Sigma}}}

% Tensor
\DeclareMathAlphabet{\mathsfit}{\encodingdefault}{\sfdefault}{m}{sl}
\SetMathAlphabet{\mathsfit}{bold}{\encodingdefault}{\sfdefault}{bx}{n}
\newcommand{\tens}[1]{\bm{\mathsfit{#1}}}
\def\tA{{\tens{A}}}
\def\tB{{\tens{B}}}
\def\tC{{\tens{C}}}
\def\tD{{\tens{D}}}
\def\tE{{\tens{E}}}
\def\tF{{\tens{F}}}
\def\tG{{\tens{G}}}
\def\tH{{\tens{H}}}
\def\tI{{\tens{I}}}
\def\tJ{{\tens{J}}}
\def\tK{{\tens{K}}}
\def\tL{{\tens{L}}}
\def\tM{{\tens{M}}}
\def\tN{{\tens{N}}}
\def\tO{{\tens{O}}}
\def\tP{{\tens{P}}}
\def\tQ{{\tens{Q}}}
\def\tR{{\tens{R}}}
\def\tS{{\tens{S}}}
\def\tT{{\tens{T}}}
\def\tU{{\tens{U}}}
\def\tV{{\tens{V}}}
\def\tW{{\tens{W}}}
\def\tX{{\tens{X}}}
\def\tY{{\tens{Y}}}
\def\tZ{{\tens{Z}}}


% Graph
\def\gA{{\mathcal{A}}}
\def\gB{{\mathcal{B}}}
\def\gC{{\mathcal{C}}}
\def\gD{{\mathcal{D}}}
\def\gE{{\mathcal{E}}}
\def\gF{{\mathcal{F}}}
\def\gG{{\mathcal{G}}}
\def\gH{{\mathcal{H}}}
\def\gI{{\mathcal{I}}}
\def\gJ{{\mathcal{J}}}
\def\gK{{\mathcal{K}}}
\def\gL{{\mathcal{L}}}
\def\gM{{\mathcal{M}}}
\def\gN{{\mathcal{N}}}
\def\gO{{\mathcal{O}}}
\def\gP{{\mathcal{P}}}
\def\gQ{{\mathcal{Q}}}
\def\gR{{\mathcal{R}}}
\def\gS{{\mathcal{S}}}
\def\gT{{\mathcal{T}}}
\def\gU{{\mathcal{U}}}
\def\gV{{\mathcal{V}}}
\def\gW{{\mathcal{W}}}
\def\gX{{\mathcal{X}}}
\def\gY{{\mathcal{Y}}}
\def\gZ{{\mathcal{Z}}}

% Sets
\def\sA{{\mathbb{A}}}
\def\sB{{\mathbb{B}}}
\def\sC{{\mathbb{C}}}
\def\sD{{\mathbb{D}}}
% Don't use a set called E, because this would be the same as our symbol
% for expectation.
\def\sF{{\mathbb{F}}}
\def\sG{{\mathbb{G}}}
\def\sH{{\mathbb{H}}}
\def\sI{{\mathbb{I}}}
\def\sJ{{\mathbb{J}}}
\def\sK{{\mathbb{K}}}
\def\sL{{\mathbb{L}}}
\def\sM{{\mathbb{M}}}
\def\sN{{\mathbb{N}}}
\def\sO{{\mathbb{O}}}
\def\sP{{\mathbb{P}}}
\def\sQ{{\mathbb{Q}}}
\def\sR{{\mathbb{R}}}
\def\sS{{\mathbb{S}}}
\def\sT{{\mathbb{T}}}
\def\sU{{\mathbb{U}}}
\def\sV{{\mathbb{V}}}
\def\sW{{\mathbb{W}}}
\def\sX{{\mathbb{X}}}
\def\sY{{\mathbb{Y}}}
\def\sZ{{\mathbb{Z}}}

% Entries of a matrix
\def\emLambda{{\Lambda}}
\def\emA{{A}}
\def\emB{{B}}
\def\emC{{C}}
\def\emD{{D}}
\def\emE{{E}}
\def\emF{{F}}
\def\emG{{G}}
\def\emH{{H}}
\def\emI{{I}}
\def\emJ{{J}}
\def\emK{{K}}
\def\emL{{L}}
\def\emM{{M}}
\def\emN{{N}}
\def\emO{{O}}
\def\emP{{P}}
\def\emQ{{Q}}
\def\emR{{R}}
\def\emS{{S}}
\def\emT{{T}}
\def\emU{{U}}
\def\emV{{V}}
\def\emW{{W}}
\def\emX{{X}}
\def\emY{{Y}}
\def\emZ{{Z}}
\def\emSigma{{\Sigma}}

% entries of a tensor
% Same font as tensor, without \bm wrapper
\newcommand{\etens}[1]{\mathsfit{#1}}
\def\etLambda{{\etens{\Lambda}}}
\def\etA{{\etens{A}}}
\def\etB{{\etens{B}}}
\def\etC{{\etens{C}}}
\def\etD{{\etens{D}}}
\def\etE{{\etens{E}}}
\def\etF{{\etens{F}}}
\def\etG{{\etens{G}}}
\def\etH{{\etens{H}}}
\def\etI{{\etens{I}}}
\def\etJ{{\etens{J}}}
\def\etK{{\etens{K}}}
\def\etL{{\etens{L}}}
\def\etM{{\etens{M}}}
\def\etN{{\etens{N}}}
\def\etO{{\etens{O}}}
\def\etP{{\etens{P}}}
\def\etQ{{\etens{Q}}}
\def\etR{{\etens{R}}}
\def\etS{{\etens{S}}}
\def\etT{{\etens{T}}}
\def\etU{{\etens{U}}}
\def\etV{{\etens{V}}}
\def\etW{{\etens{W}}}
\def\etX{{\etens{X}}}
\def\etY{{\etens{Y}}}
\def\etZ{{\etens{Z}}}

% The true underlying data generating distribution
\newcommand{\pdata}{p_{\rm{data}}}
% The empirical distribution defined by the training set
\newcommand{\ptrain}{\hat{p}_{\rm{data}}}
\newcommand{\Ptrain}{\hat{P}_{\rm{data}}}
% The model distribution
\newcommand{\pmodel}{p_{\rm{model}}}
\newcommand{\Pmodel}{P_{\rm{model}}}
\newcommand{\ptildemodel}{\tilde{p}_{\rm{model}}}
% Stochastic autoencoder distributions
\newcommand{\pencode}{p_{\rm{encoder}}}
\newcommand{\pdecode}{p_{\rm{decoder}}}
\newcommand{\precons}{p_{\rm{reconstruct}}}

\newcommand{\laplace}{\mathrm{Laplace}} % Laplace distribution

\newcommand{\E}{\mathbb{E}}
\newcommand{\Ls}{\mathcal{L}}
\newcommand{\R}{\mathbb{R}}
\newcommand{\emp}{\tilde{p}}
\newcommand{\lr}{\alpha}
\newcommand{\reg}{\lambda}
\newcommand{\rect}{\mathrm{rectifier}}
\newcommand{\softmax}{\mathrm{softmax}}
\newcommand{\sigmoid}{\sigma}
\newcommand{\softplus}{\zeta}
\newcommand{\KL}{D_{\mathrm{KL}}}
\newcommand{\Var}{\mathrm{Var}}
\newcommand{\standarderror}{\mathrm{SE}}
\newcommand{\Cov}{\mathrm{Cov}}
% Wolfram Mathworld says $L^2$ is for function spaces and $\ell^2$ is for vectors
% But then they seem to use $L^2$ for vectors throughout the site, and so does
% wikipedia.
\newcommand{\normlzero}{L^0}
\newcommand{\normlone}{L^1}
\newcommand{\normltwo}{L^2}
\newcommand{\normlp}{L^p}
\newcommand{\normmax}{L^\infty}

\newcommand{\parents}{Pa} % See usage in notation.tex. Chosen to match Daphne's book.

\DeclareMathOperator*{\argmax}{arg\,max}
\DeclareMathOperator*{\argmin}{arg\,min}

\DeclareMathOperator{\sign}{sign}
\DeclareMathOperator{\Tr}{Tr}
\let\ab\allowbreak


\usepackage{subcaption}
\usepackage{wrapfig}
\usepackage{hyperref}            % link colours are set by preprint.sty
\usepackage{url}
\usepackage[capitalise,noabbrev,nameinlink]{cleveref}
\usepackage{algorithm}
\usepackage{algpseudocode}
\usepackage{amssymb}
\usepackage{mathtools}            % \mathclap for wide \underbrace labels
\usepackage{multirow}
\usepackage{makecell}
\usepackage{float}

% ------------------------------------------------------------------ entities
% One macro per name; the font follows the entity's category, not its role
% in a table. All names below are placeholders.
\newcommand{\ourmethod}{\rpsMethod{OurMethod}}        % our method: bold mono
\newcommand{\priorA}{\rpsBaseline{PriorA}}            % external methods: plain
\newcommand{\priorB}{\rpsBaseline{PriorB}}
\newcommand{\modelS}{\rpsModel{Model-3B}}             % models: pcr
\newcommand{\modelL}{\rpsModel{Model-7B}}
\newcommand{\corpus}{\rpsDataset{TrainCorpus}}        % datasets: plain
\newcommand{\benchA}{\rpsBenchmark{BenchA}}           % benchmarks: cmtt
\newcommand{\benchB}{\rpsBenchmark{BenchB}}
% Named modules: small caps + a fixed symbol, reused in text, figures, tables.
\newcommand{\router}{\rpsComponent{95}{Router}}       % filled flower
\newcommand{\corrector}{\rpsComponent{96}{Corrector}} % open flower
\newcommand{\routername}{\rpsModule{Router}}
\newcommand{\correctorname}{\rpsModule{Corrector}}

% ------------------------------------------------------------------ front matter
\title{\ourmethod{}: Your Paper Title Goes Here\\in Two Balanced Lines}

% One \author per person; the optional argument holds the superscripts.
\author[1,2,\star]{First Author}
\author[1,\star]{Second Author}
\author[2]{Third Author}
\author[3]{Fourth Author}
\author[1,\dagger]{Last Author}

\affiliation[1]{University One}
\affiliation[2]{Company Two}
\affiliation[3]{Institute Three}

\authornote[\star]{Equal contribution}
\authornote[\dagger]{Corresponding author}

% Lines printed at the bottom of the abstract box (all optional).
\date{\today}
\metadata[Code]{\url{https://github.com/your-org/your-repo}}
\correspondence{\email{last.author@example.com}}

% Logos at the top left of page 1 (delete \logos{...} for none).
% NTU lockup | Midori icon + wordmark, all centred on one axis.
\logos{%
  \logoimage[0.42in]{figures/ntu-logo}%
  \logodivider[0.42in]%
  \logoimage[0.42in]{figures/midori-icon}\hspace{0.6em}%
  \logoimage[0.13in]{figures/midori-wordmark}}
% Alternative: the original stacked Midori logo (icon above wordmark):
%   \logoimage[0.55in]{figures/midori-logo}

% Header text from page 2 onward.
\runninghead{Preprint}

\begin{document}
\maketitle

\begin{abstract}
Template text: every claim, name and number in this file is a placeholder.
Large models often fail when \rpsStrong{two input sources disagree}, even
though the evidence needed for the answer is still present in one of them.
We propose \ourmethod{}, a \rpsTerm{training-free} method that detects such
disagreements and corrects the joint prediction at inference time. On \benchA{}
and \benchB{}, \ourmethod{} improves overall accuracy on \modelL{} from 70.5\%
to 85.9\% in this illustrative setting.
\end{abstract}

\section{Introduction}
\label{sec:intro}

Background paragraphs stay in plain text. Recent models jointly process
several input sources and are widely used in downstream
applications~\citep{vaswani2017attention,goodfellow2016deep}. Combining sources,
however, does not guarantee that the answer is grounded in the right one.

The key bottleneck is that \rpsStrong{the relevant evidence is often still
recoverable}, yet the joint prediction ignores it. The two sources differ in
\rpsTerm{which question they can answer}, while sharing the same
\rpsTerm{internal representation}. This raises the central question of this work:

\begin{rpsquestion}
Can a model recover the answer from the relevant source when the joint
prediction has already gone wrong?
\end{rpsquestion}

We answer it with \ourmethod{} (\cref{fig:overview}), which first locates the
relevant source with a \router{} and then steers the joint representation
with a \corrector{}. The main contributions are:

\begin{contributions}
  \item \rpsStrongTerm{A recoverable failure mode.} We identify cases where the
  joint prediction is wrong but a single-source view is right, and measure how
  often they occur.
  \item \rpsStrongTerm{A training-free correction.} \ourmethod{} detects these
  cases and corrects them at inference time without updating model weights.
  \item \rpsStrongTerm{Controlled evaluation.} We evaluate on \benchA{} and
  \benchB{} with two model sizes and compare against \priorA{} and \priorB{}.
\end{contributions}

\begin{figure}[t]
  \centering
  \includegraphics[width=\textwidth,height=1.8in]{example-image-16x9}
  \caption{\textbf{Overview of \ourmethod{}.} \rpsStrongTerm{(a) Detect:} the
  \router{} selects the relevant source and checks for a decision conflict.
  \rpsStrongTerm{(b) Correct:} the \corrector{} steers the joint representation
  toward that source. Flower markers match the module names in the text.}
  \label{fig:overview}
\end{figure}

\section{Related Work}
\label{sec:related}

\paragraph{Joint multi-source models.}
Plain text for the survey of one category of prior work~\citep{Hinton06}.
End with one sentence that positions this paper against the category.

\paragraph{Inference-time correction.}
Methods such as \priorA{} and \priorB{} adjust the output distribution at
decoding time. They \rpsTerm{suppress} errors, whereas \ourmethod{}
\rpsTerm{recovers} the missing evidence.

\section{Method}
\label{sec:method}

\paragraph{Problem setup.}
Given input sources $a$ and $v$ and a question $q$, a model with parameters
$\bm{\theta}$ produces a score $z_c$ for every candidate $c\in\mathcal{C}$.
The prediction and its margin are
\begin{equation}
  \hat{y} = \argmax_{c\in\mathcal{C}} z_c, \qquad
  m = z_{\hat{y}} - \max_{c\neq\hat{y}} z_c .
  \label{eq:margin}
\end{equation}

\noindent\rpsStage{Stage 1: Detect.}
The \router{} withholds the competing source and predicts $\hat{y}_{\mathrm{uni}}$.
A conflict is flagged when $\hat{y}_{\mathrm{uni}}\neq\hat{y}$.

\noindent\rpsStage{Stage 2: Correct.}
For flagged inputs, the \corrector{} shifts the hidden state
$\mathbf{h}_\ell\in\mathbb{R}^{d}$ at layer $\ell$:
\begin{equation}
  \mathbf{h}'_\ell
  = \mathbf{h}_\ell
  + \underbrace{\alpha\,(\mathbf{h}^{\mathrm{uni}}_\ell - \mathbf{h}_\ell)}_{\mathclap{\text{\correctorname{} residual}}},
  \label{eq:correct}
\end{equation}
where $\alpha$ scales with the confidence of the single-source view.
\cref{alg:method} summarizes both stages.

\begin{algorithm}[t]
\caption{\ourmethod{} inference}
\label{alg:method}
\begin{algorithmic}[1]
\Require sources $a, v$; question $q$; strength $\alpha$
\State $\hat{y} \gets \argmax_c z_c(a, v, q)$ \Comment{joint prediction}
\State $\hat{y}_{\mathrm{uni}} \gets$ \routername{}$(a, v, q)$ \Comment{Stage 1: Detect}
\If{$\hat{y}_{\mathrm{uni}} \neq \hat{y}$}
  \State apply \cref{eq:correct} at the selected layers \Comment{Stage 2: Correct}
  \State $\hat{y} \gets \argmax_c z'_c$
\EndIf
\State \Return $\hat{y}$
\end{algorithmic}
\end{algorithm}

\section{Experiments}
\label{sec:exp}

\paragraph{Setup.}
We evaluate \modelS{} and \modelL{} on \benchA{} and \benchB{}; the
hyperparameters are tuned on a held-out split of \corpus{}. Baselines are
\priorA{} and \priorB{} under the same decoding settings.

\begin{table*}[t]
\centering
\caption{\textbf{Main results (\%).} Placeholder numbers on \modelL{}. Higher
is better for all columns. Within each column, the best value is \textbf{bold}
and the second best is \underline{underlined}; the shaded row marks our method
and does not by itself indicate the best result.}
\label{tab:main}
\small
\setlength{\tabcolsep}{4pt}
\renewcommand{\arraystretch}{1.12}
\begin{tabular}{@{}ll cc cc c@{}}
\toprule
& & \multicolumn{2}{c}{\benchA{}} & \multicolumn{2}{c}{\benchB{}}
& \multirow{2}{*}{\shortstack[c]{Overall\\Acc.}} \\
\cmidrule(lr){3-4}\cmidrule(lr){5-6}
\textbf{Backbone} & \textbf{Method} & Acc. & F1 & Acc. & F1 & \\
\midrule
\multirow{4}{*}{\modelL{}}
& Baseline   & 69.3 & 75.7 & 71.8 & 77.0 & 70.5 \\
& +\priorA{} & \rpsSecond{75.5} & \rpsSecond{79.4} & 66.5 & 73.5 & 71.0 \\
& +\priorB{} & 72.1 & 77.8 & \rpsSecond{75.5} & \rpsSecond{79.1} & \rpsSecond{73.8} \\
& \ourscell{+\ourmethod{}}
& \ourscell{\rpsBest{82.0}} & \ourscell{\rpsBest{84.1}}
& \ourscell{\rpsBest{89.8}} & \ourscell{\rpsBest{89.8}}
& \ourscell{\rpsBest{85.9}} \\
\bottomrule
\end{tabular}
\end{table*}

\rpsObservation{\ourmethod{} gives the largest gain on \benchB{}.}\label{obs:bench}
Relative to the strongest baseline in each column, accuracy rises by 6.5
points on \benchA{} (75.5 $\rightarrow$ 82.0) and by 14.3 points on \benchB{}
(75.5 $\rightarrow$ 89.8), as shown in \cref{tab:main}.

\begin{wraptable}{r}{0.46\columnwidth}
\vspace{-0.6em}
\setlength{\abovecaptionskip}{3pt}
\setlength{\belowcaptionskip}{-6pt}
\setlength{\tabcolsep}{2.7pt}
\renewcommand{\arraystretch}{1.05}
\centering
\caption{\textbf{Scaling.} Overall accuracy (\%); gains over Base in parentheses.}
\label{tab:scaling}
\small
\begin{tabular}{llc}
\toprule
\textbf{Model} & \textbf{Method} & \textbf{Acc.} \\
\midrule
\multirow{2}{*}{\modelS{}} & Base & 76.9 \\
 & $\hookrightarrow$ +\ourmethod{} & 82.5 \gain{5.6} \\
\midrule
\multirow{2}{*}{\modelL{}} & Base & 70.5 \\
 & $\hookrightarrow$ +\ourmethod{} & 85.9 \gain{15.4} \\
\bottomrule
\end{tabular}
\end{wraptable}

\rpsObservation{Across the two sizes we test, the gain is larger on \modelL{}.}\label{obs:scale}
\cref{tab:scaling} shows a gain of 5.6 points on \modelS{} and 15.4 points on
\modelL{}. With only two sizes, this does not establish a trend with scale.
Plain text continues the analysis here. Lorem ipsum dolor sit amet,
consectetur adipiscing elit, sed do eiusmod tempor incididunt ut labore et
dolore magna aliqua. Ut enim ad minim veniam, quis nostrud exercitation
ullamco laboris nisi ut aliquip ex ea commodo consequat.

\begin{rpsinsight}{Takeaway}
\rpsStrong{The overall gain is not uniform across benchmarks or model sizes.}
Report \benchA{} and \benchB{} separately, as in \cref{obs:bench,obs:scale}.
\end{rpsinsight}

\section{Conclusion}
Summarize the findings in plain text, reusing the entity macros.

\paragraph{Limitations and Future Work.}
\ourmethod{} assumes that \rpsStrong{the relevant evidence is still present in
one source}; it may help less when that evidence is weakly encoded.

\subsection*{AI use statement}
Describe how generative AI tools were used in preparing the manuscript.

\subsection*{Reproducibility Statement}
Point to the sections and appendices that support reproduction.

\bibliography{references}
\bibliographystyle{preprint}

% ------------------------------------------------------------------ appendix
\appendix
\clearpage

\section*{Appendix Overview}

\begin{appendixoverview}
  \appentry{app:impl}{Implementation Details}
  \appsubentry{app:prompt}{Prompt for the \routername{}}
  \appendrule
  \appentry{app:style}{Typography Key}
\end{appendixoverview}

\section{Implementation Details}
\label{app:impl}

\subsection{Prompt for the \routername{}}
\label{app:prompt}
Prompts and other raw materials go in a breakable monospace box; the text is
kept verbatim.

\begin{rpsprompt}{Router prompt}
SYSTEM: You will see a question about an audio-visual clip.
USER: {question}
To answer this question, which modality is needed (audio, video, or both):
\end{rpsprompt}

\section{Typography Key}
\label{app:style}

\begin{table}[H]
\centering
\caption{\textbf{Formatting roles used in this template.} Delete this appendix
in your paper.}
\label{tab:style-key}
\small
\begin{tabular}{lll}
\toprule
\textbf{Role} & \textbf{Macro} & \textbf{Example} \\
\midrule
Our method        & \rpsCode{\char92 rpsMethod}    & \ourmethod{} \\
External method   & \rpsCode{\char92 rpsBaseline}  & \priorA{} \\
Model             & \rpsCode{\char92 rpsModel}     & \modelL{} \\
Dataset           & \rpsCode{\char92 rpsDataset}   & \corpus{} \\
Benchmark         & \rpsCode{\char92 rpsBenchmark} & \benchA{} \\
Named module      & \rpsCode{\char92 rpsComponent} & \router{}, \corrector{} \\
Method stage      & \rpsCode{\char92 rpsStage}     & \rpsStage{Stage 1: Detect.} \\
Key phrase        & \rpsCode{\char92 rpsStrong}    & \rpsStrong{bottleneck} \\
Contrast          & \rpsCode{\char92 rpsTerm}      & \rpsTerm{suppress} vs.\ \rpsTerm{recover} \\
Contribution head & \rpsCode{\char92 rpsStrongTerm} & \rpsStrongTerm{Short title.} \\
Finding           & \rpsCode{\char92 rpsObservation} & \rpsStrong{Observation \rpsNumber{1}:} \ldots \\
Gain / drop       & \rpsCode{\char92 gain}, \rpsCode{\char92 drop} & 85.9 \gain{15.4}, 70.1 \drop{2.3} \\
\bottomrule
\end{tabular}
\end{table}

\end{document}
